% !TeX program = lualatex
% Compile twice: lualatex 80.tex
% Standalone research notebook; original section preserved.
% Research additions: 27 September 2026. No external TeX input or BibTeX file.
\documentclass[11pt,a4paper]{article}
\usepackage[margin=24mm,headheight=14pt]{geometry}
\usepackage{fontspec}
\setmainfont{Latin Modern Roman}
\setsansfont{Latin Modern Sans}
\setmonofont{Latin Modern Mono}
\usepackage{amsmath,amssymb,mathtools}
\usepackage{unicode-math}
\setmathfont{Latin Modern Math}
\usepackage{microtype}
\usepackage{enumitem,xurl,needspace}
\usepackage[dvipsnames]{xcolor}
\usepackage{hyperref}
\hypersetup{unicode=true,colorlinks=true,linkcolor=MidnightBlue,urlcolor=MidnightBlue,
 pdftitle={Research notebook - Problem 80},pdfauthor={Open Problems Math},bookmarksnumbered=true}
\usepackage{bookmark}
\usepackage{titlesec}
\titleformat{\section}{\Large\bfseries\raggedright}{\thesection}{0.7em}{}
\usepackage{fancyhdr}
\pagestyle{fancy}\fancyhf{}
\fancyhead[L]{\small\sffamily Open Problems Math | Research notebook}
\fancyhead[R]{\small\sffamily Problem 80 | 27 September 2026}
\fancyfoot[C]{\thepage}
\setlength{\parindent}{0pt}
\setlength{\parskip}{5pt plus 1pt}
\setlength{\emergencystretch}{3em}
\setcounter{tocdepth}{1}
\setcounter{secnumdepth}{1}
\setlist[itemize]{leftmargin=3em,itemsep=5pt,topsep=4pt,parsep=0pt}
\allowdisplaybreaks
\numberwithin{equation}{section}
\newcommand{\N}{\mathbb{N}}
\newcommand{\Z}{\mathbb{Z}}
\newcommand{\Q}{\mathbb{Q}}
\newcommand{\R}{\mathbb{R}}
\newcommand{\C}{\mathbb{C}}
\newcommand{\F}{\mathbb{F}}
\newcommand{\A}{\mathbb{A}}
\newcommand{\PP}{\mathbb P}
\newcommand{\E}{\mathbb{E}}
\newcommand{\eps}{\varepsilon}
\newcommand{\dd}{\,\mathrm d}
\newcommand{\sref}[2]{\hyperlink{source:#1:#2}{[S#2]}}
\newcommand{\eref}[2]{\hyperlink{extra:#1:#2}{[E#2]}}
\newif\ifshowcatalogue
\showcataloguetrue
\newcommand{\cataloguescope}[1]{\ifshowcatalogue
 \paragraph{Exact scope in the supplied catalogue.}
 \begin{quote}\small #1\end{quote}\fi}
\begin{document}
\setcounter{section}{79}
\section{\texorpdfstring{Hadwiger Conjecture}{Hadwiger Conjecture}}\label{80}
\subsection{Definitions and mathematical statement}
A graph minor is obtained by vertex deletions, edge deletions, and edge contractions. A proper $t$-coloring is a map $c:V(G)\to\{1,\ldots,t\}$ with different colors at adjacent vertices. Hadwiger's assertion is
\[
 K_{t+1}\not\preccurlyeq G\quad\Longrightarrow\quad\chi(G)\le t
\]
for every finite simple graph $G$ and every $t\ge1$. The complete graph is excluded as a \emph{minor}, not merely as a subgraph or an induced subgraph.
\par\smallskip\textit{Formulation and concept references:} \sref{80}{1}.
\cataloguescope{Prove that every simple graph with no K\_\{t+1\} minor is t-colourable for every integer t >= 1.}
\subsection{Short English statement}
If a graph cannot contain a complete graph of order $t+1$ even after contractions, must $t$ colors suffice?
% BEGIN RESEARCH_ADDITION ATTEMPT_80
\subsection{Research attempt}

\paragraph{Current frontier.}
Delcourt--Postle's established asymptotic colouring bound is
$O(t\log\log t)$ for graphs excluding $K_t$ as a minor. Liu--Luo's September
2026 preprint states an improvement to $O(t\log\log\log t)$; this is a
recent authors' theorem, not independently proof-audited here
\eref{80}{1}. Neither bound gives the exact $t-1$ required for $K_t$-minor-free
graphs. Claims about the odd-minor variant in \sref{80}{3} do not by
themselves refute the ordinary minor assertion. The investigation below
uses ordinary connected branch sets throughout.

\paragraph{Attack route.}
A minimum counterexample would be vertex-critical and have strong degree
constraints, but large degree alone does not force the required exact
clique minor. Another route contracts a connected dominating set to create
one branch set and inducts on the remaining graph. This gives an elementary
known independence-number bound. The missing step for Hadwiger is then
pinpointed as retention of enough chromatic number after removing such a
set. The proof is worked out rather than presenting the weaker bound as
the target.

\paragraph{Attempt: a connected dominating set with an independent witness.}
For a connected nonempty graph $G$, write $\alpha=\alpha(G)$. Construct
sets $I\subseteq D$ with $I$ independent, $D$ connected, and
$|D|=2|I|-1$. Start with one vertex in both sets. If $D$ does not dominate
$G$, connectedness supplies a vertex $y$ at distance exactly two from
$D$, on a path $d,x,y$ with $d\in D$. Add $x,y$ to $D$ and $y$ to $I$.
The new $D$ is connected, and $y$ has no neighbour in the old $D$, so the
new $I$ stays independent. The stated size invariant persists. Since the
graph is finite, the process stops at a connected dominating set with
\begin{equation}\label{eq80:dom}
 |D|\leq2\alpha(G)-1.
\end{equation}
This is the elementary connected-domination mechanism behind the classical
Duchet--Meyniel bound \eref{80}{2}; it is reconstructed here without a
novelty claim.

Let $h(G)$ be the maximum order of a complete minor. We now prove
\begin{equation}\label{eq80:weak}
 h(G)\geq\frac{|V(G)|}{2\alpha(G)-1}
\end{equation}
for every nonempty graph, by induction on its number of vertices.
For disconnected graphs with nonempty components $G_j$, induction gives
$h(G_j)\geq n_j/(2\alpha_j-1)$. Since
$\alpha(G)=\sum_j\alpha_j$ and $h(G)=\max_jh(G_j)$,
\[
 h(G)\geq\frac{\sum_j n_j}{\sum_j(2\alpha_j-1)}
       \geq\frac{|V(G)|}{2\alpha(G)-1}.
\]
For connected $G$ choose $D$ from \eqref{eq80:dom}. If $D=V(G)$, the
right side of \eqref{eq80:weak} is at most one. Otherwise let $G_j$ be
the components of $G-D$. Contract $D$ to one branch set. It is adjacent
to every vertex outside it, hence to every branch set of any clique minor
in any $G_j$. Thus
\[
 h(G)\geq1+\max_jh(G_j)
 \geq1+\frac{|V(G)|-|D|}{\sum_j(2\alpha(G_j)-1)}
 \geq1+\frac{|V(G)|-|D|}{2\alpha(G)-1}.
\]
The last quantity is at least the right side of \eqref{eq80:weak}, by
\eqref{eq80:dom}. This completes the induction, with the empty remainder
handled separately rather than dividing by an empty denominator.

\paragraph{Attempt: a precise chromatic-retention condition.}
Suppose every connected graph of chromatic number $k\geq2$ possessed a
connected dominating set $D$ with
\begin{equation}\label{eq80:missing}
 \chi(G-D)\geq k-1.
\end{equation}
Then induction on the number of vertices would prove $h(G)\geq\chi(G)$:
the smaller graph $G-D$ would have a $K_{k-1}$ minor, and contracting $D$
would supply the $k$th branch set. Disconnected graphs are handled by a
component of maximum chromatic number. Thus \eqref{eq80:missing} is one
explicit sufficient lemma for the full assertion, not something implied
by \eqref{eq80:dom}.

The converse direction need not be asserted. Hadwiger might hold by branch
sets that cannot be arranged through this dominating-set induction. The
condition is a proposed stronger structural route, and a counterexample
to it would not itself refute Hadwiger.

There is nevertheless a rigorous minimum-counterexample consequence.
If a smallest-vertex graph $G$ violates $h(G)\geq\chi(G)$, put
$k=\chi(G)$. It is connected and vertex-critical: otherwise a smaller
component or a vertex deletion has the same chromatic number, contradicting
minimality. Hence its minimum degree is at least $k-1$, because a colouring
of $G-v$ by $k-1$ colours would extend across a vertex of degree at most
$k-2$. Moreover, \emph{every} connected dominating set $D$ in this $G$
satisfies $\chi(G-D)\leq k-2$. Otherwise minimality and the contraction
argument just given contradict $h(G)<k$. This is a concrete constraint a
counterexample search can test, with no assumption that such a graph exists.

\paragraph{Stress test.}
A bound on $|D|$ in terms of $\alpha(G)$ does not control the loss of
chromatic number. The elementary inequality
$\chi(G)\leq\chi(G[D])+\chi(G-D)$ loses as many colours as the induced
subgraph on $D$ needs; connectedness does not make that number one.
The independent witness $I\subset D$ constructed above certifies the size
bound, not a colouring of all of $D$ with one colour. This is precisely the
unproved leap that would turn \eqref{eq80:weak} into Hadwiger.

Kempe-chain ideas face a related issue: paths between colour classes may
intersect each other. A clique minor needs disjoint connected branch sets,
not merely a path for each desired pair. The connected dominating set
solves disjointness for one branch set, but the retained chromatic number
is then the obstacle. Large average degree, restricted forbidden subgraphs,
or a bound with extra logarithmic factors do not supply the exact condition.

The companion script tests the construction \eqref{eq80:dom} on a finite
collection of small connected graphs and verifies connectedness, domination,
independence and the size identity after every step. It does not certify
\eqref{eq80:missing} for all graphs or search all possible minors at
unbounded order. The known bound has an elementary proof here, so numerical
agreement is not being used as its justification.

\paragraph{Outcome.}
\textbf{Outcome: Exploratory approach with unresolved gap.}

The attempt reconstructs a rigorous connected-domination lower bound for
clique minors and isolates a sufficient chromatic-retention lemma. It also
proves what any smallest counterexample would have to do to every connected
dominating set. No exact colouring theorem beyond the established cases
is obtained, and no counterexample to the original conjecture is produced.
A continuation needs new control of chromatic loss or a different disjoint
branch-set construction; the size estimate alone cannot provide it.
% END RESEARCH_ADDITION ATTEMPT_80
\subsection{Sources}
\begin{itemize}\raggedright
\item[\textnormal{[S1]}]\hypertarget{source:80:1}{} Lafferty and Song, Every graph with no K\_8\textasciicircum{}\{-4\} minor is 7-colorable (2022), abstract.
\url{https://arxiv.org/abs/2208.07338}.
{\small\textit{Catalogue formulation source.}}
\item[\textnormal{[S2]}]\hypertarget{source:80:2}{} Beyond halfway to Hadwiger's conjecture.
\url{https://arxiv.org/abs/2609.06867v2}.
{\small\textit{Catalogue research/status reference; not a proof endorsement.}}
\item[\textnormal{[S3]}]\hypertarget{source:80:3}{} Disproof of the Odd Hadwiger Conjecture.
\url{https://arxiv.org/abs/2512.20392}.
{\small\textit{Catalogue research/status reference; not a proof endorsement.}}
% BEGIN RESEARCH_ADDITION REFERENCES_80
\item[\textnormal{[E1]}]\hypertarget{extra:80:1}{} Chun-Hung Liu and
Jason Luo, \emph{Beyond halfway to Hadwiger's conjecture},
arXiv:2609.06867v2, 9 September 2026, abstract and main theorem.
\url{https://arxiv.org/abs/2609.06867v2}.
{\small\textit{Inspected 27 September 2026; the $O(t\log\log\log t)$
claim is a recent preprint theorem and not independently audited here.
The prior Delcourt--Postle bound is described in its introduction.}}
\item[\textnormal{[E2]}]\hypertarget{extra:80:2}{} Pierre Duchet and
Henri Meyniel, \emph{On Hadwiger's number and the stability number},
in \emph{Graph Theory}, North-Holland Mathematics Studies 62 (1982),
71--73, DOI 10.1016/S0304-0208(08)73549-7.
\url{https://doi.org/10.1016/S0304-0208(08)73549-7}.
{\small\textit{Classical attribution; the connected-domination proof and
all inequalities used in this attempt are supplied explicitly.}}
% END RESEARCH_ADDITION REFERENCES_80
\end{itemize}

\end{document}
